% PAH — PM01: 6 CAU DUNG/SAI KHONG SO, BAN THAO
% LaTeX BODY su dung ex_test: \begin{ex}, \choiceTF, \loigiai
% Nguon tham chieu: xem chuyen_de/PM01_TIEN_TRINH_TU_DUY_KHONG_SO_V1.md

% PM01 Q1: Nguồn điện cấp nhiệt có phải chất nhận hết? (Mức 1)
\begin{ex}
Một phòng thực hành muốn xác định nhiệt nóng chảy riêng của một chất rắn tinh thể. Chất được duy trì ở nhiệt độ nóng chảy, cao hơn nhiệt độ phòng, trong suốt hai giai đoạn. Người làm thí nghiệm dùng cùng một nồi đun, mỗi giai đoạn cấp điện trong thời gian $t$, nhưng công suất lần sau $P_2$ lớn hơn $P_1$. Phần nhiệt truyền ra môi trường trong hai giai đoạn được coi là bằng nhau và khác không; còn đủ chất rắn chưa nóng chảy hết ở cả hai giai đoạn. Gọi $m_1,m_2$ là khối lượng nóng chảy tương ứng.
\choiceTF
{\True Điện năng dây đốt cung cấp ở giai đoạn $i$ là $P_it$.}
{Nhiệt lượng chất rắn dùng để nóng chảy ở giai đoạn $i$ bằng toàn bộ $P_it$.}
{\True Nếu những điều kiện trên được bảo đảm, khối lượng nóng chảy ở giai đoạn 2 lớn hơn giai đoạn 1.}
{\True Việc giữ thời gian bằng nhau không cản trở xác định nhiệt nóng chảy riêng từ hai kết quả đo khác nhau.}
\loigiai{Đáp án: a) Đúng; b) Sai; c) Đúng; d) Đúng. Xem lời giải và điều kiện vật lí tại hồ sơ PM01 trên Git.}
\end{ex}

% PM01 Q2: Khi khối lượng gấp đôi, nhiệt độ có giảm một nửa? (Mức 2)
\begin{ex}
Một nhóm học sinh muốn đo nhiệt dung riêng của một chất lỏng trong cùng nhiệt lượng kế. Một bạn dự đoán: ``Giữ nguyên điện năng đun, khi khối lượng chất lỏng gấp đôi thì độ tăng nhiệt độ giảm còn đúng một nửa.'' Nhóm đun hai khối lượng $m$ và $2m$, cùng điện năng, cùng nhiệt độ ban đầu và đo hai độ tăng nhiệt độ. Chất không chuyển thể; nhiệt dung hiệu dụng $C$ của bình dương, không đổi; bỏ qua trao đổi nhiệt với bên ngoài và các ảnh hưởng khác.
\choiceTF
{\True Điện năng cấp cho mỗi lần đo bằng tổng nhiệt lượng chất lỏng và nhiệt lượng kế hấp thụ.}
{Vì khối lượng lần sau gấp đôi nên độ tăng nhiệt độ lần sau bằng đúng một nửa lần đầu.}
{\True Độ tăng nhiệt độ lần sau nhỏ hơn lần đầu nhưng lớn hơn một nửa lần đầu.}
{Nếu giữ khối lượng không đổi mà chỉ thay công suất và thời gian, hai lần đo vẫn luôn cho phép tìm riêng cả nhiệt dung riêng của chất và nhiệt dung của bình.}
\loigiai{Đáp án: a) Đúng; b) Sai; c) Đúng; d) Sai. Xem lời giải và điều kiện vật lí tại hồ sơ PM01 trên Git.}
\end{ex}

% PM01 Q3: Vai trò của môi trường khi chất lỏng đang sôi (Mức 2)
\begin{ex}
Một chất lỏng tinh khiết có nhiệt độ sôi thấp hơn nhiệt độ phòng được khảo sát bằng nguồn điện điều chỉnh công suất. Khi chất đã sôi ổn định, người nghiên cứu cấp điện hai giai đoạn cùng thời gian $t$, với $P_2>P_1\ge0$, đo khối lượng hóa hơi $m_1,m_2$. Coi nhiệt lượng dương từ môi trường truyền vào chất lỏng ở mỗi giai đoạn như nhau, nhiệt hóa hơi riêng không đổi, bỏ qua các ảnh hưởng khác.
\choiceTF
{Nguồn điện là nguồn năng lượng duy nhất làm chất lỏng hóa hơi.}
{\True Hiệu nhiệt lượng dùng để hóa hơi ở hai giai đoạn bằng $(P_2-P_1)t$.}
{\True Nếu ở một giai đoạn ngừng cấp điện, chất lỏng vẫn có thể tiếp tục hóa hơi nhờ nhiệt lượng môi trường truyền vào, trong các điều kiện đang xét.}
{Khối lượng hóa hơi luôn tỉ lệ thuận với công suất điện cấp, tức $m_2/m_1=P_2/P_1$ với $P_1>0$.}
\loigiai{Đáp án: a) Sai; b) Đúng; c) Đúng; d) Sai. Xem lời giải và điều kiện vật lí tại hồ sơ PM01 trên Git.}
\end{ex}

% PM01 Q4: Hai thời gian khác nhau có làm hỏng phương pháp? (Mức 3)
\begin{ex}
Một thiết bị đang đo nhiệt hóa hơi riêng của chất lỏng có nhiệt độ sôi cao hơn nhiệt độ phòng. Kỹ thuật viên thực hiện hai giai đoạn ở hai công suất và hai khoảng thời gian có thể khác nhau. Từ công suất, thời gian và sự giảm khối lượng, kỹ thuật viên xác định $P_i,t_i,m_i$ ($t_i>0$). Giả sử công suất nhiệt mất ra môi trường là một hằng số dương như nhau trong hai giai đoạn và nhiệt hóa hơi riêng không đổi.
\choiceTF
{\True Nhiệt lượng mất ra ngoài ở mỗi giai đoạn tỉ lệ với thời gian giai đoạn ấy.}
{\True Có thể chia phương trình năng lượng của từng giai đoạn cho thời gian rồi lấy hiệu để loại công suất mất nhiệt.}
{Nếu hai thời gian khác nhau, chỉ cần lấy hiệu hai điện năng $P_2t_2-P_1t_1$ là luôn bằng hiệu nhiệt lượng dùng hóa hơi $(m_2-m_1)L$.}
{Nếu hai phép đo có cùng tốc độ hóa hơi $m_1/t_1=m_2/t_2$, dữ liệu vẫn luôn xác định duy nhất được cả nhiệt hóa hơi riêng và công suất mất nhiệt.}
\loigiai{Đáp án: a) Đúng; b) Đúng; c) Sai; d) Sai. Xem lời giải và điều kiện vật lí tại hồ sơ PM01 trên Git.}
\end{ex}

% PM01 Q5: Hai cách bố trí, cách nào thực sự tách được hai ẩn? (Mức 3)
\begin{ex}
Hai nhóm cùng xác định nhiệt dung riêng của một chất lỏng bằng bộ nhiệt lượng kế như nhau. Nhiệt dung hiệu dụng của bình $C$ chưa biết. Nhóm A giữ nguyên khối lượng chất lỏng, thay đổi công suất và thời gian rồi đo độ tăng nhiệt độ. Nhóm B đổi khối lượng chất lỏng giữa hai lần đo và cũng ghi đủ điện năng cùng độ tăng nhiệt độ. Trong mọi phép đo, các độ tăng nhiệt độ đều dương; chất không chuyển thể; $c,C$ không đổi và không có thất thoát nhiệt ra môi trường.
\choiceTF
{\True Dữ liệu nhóm A có thể xác định tổ hợp $mc+C$.}
{Nếu hai lần đo của nhóm A có độ tăng nhiệt độ khác nhau, điều đó đủ để suy ra riêng $c$ và $C$.}
{\True Với hai khối lượng khác nhau, nhóm B có thể xác định riêng $c$ và $C$ trong mô hình đã nêu.}
{Thêm một lần đo thứ ba với đúng khối lượng ban đầu luôn giúp nhóm A xác định riêng $c$ và $C$.}
\loigiai{Đáp án: a) Đúng; b) Sai; c) Đúng; d) Sai. Xem lời giải và điều kiện vật lí tại hồ sơ PM01 trên Git.}
\end{ex}

% PM01 Q6: Khi kết quả đo không phù hợp với mô hình (Mức 4 — ĐÍCH)
\begin{ex}
Trong một nghiên cứu xác định nhiệt hóa hơi riêng, một chất lỏng được giữ ở nhiệt độ sôi lớn hơn nhiệt độ phòng. Nhà nghiên cứu giả định nhiệt hóa hơi riêng không đổi và công suất thất thoát nhiệt ra môi trường là như nhau ở mọi giai đoạn. Sau khi tiến hành hai giai đoạn với công suất cấp điện ở giai đoạn 2 lớn hơn giai đoạn 1, người nghiên cứu ghi nhận khối lượng hóa hơi trên mỗi đơn vị thời gian lại bằng nhau giữa hai giai đoạn.  Một thành viên kết luận rằng phép đo đã xác nhận mô hình ban đầu. Các thành viên khác đề xuất kiểm tra lại quy trình và giả thiết.
\choiceTF
{\True Theo mô hình đã giả định, công suất cấp lớn hơn phải tương ứng với tốc độ hóa hơi lớn hơn.}
{Kết quả vừa ghi nhận là bằng chứng đủ để xác nhận công suất thất thoát nhiệt không đổi giữa hai giai đoạn.}
{Kết quả trái dự đoán chứng minh rằng định luật bảo toàn năng lượng không còn đúng trong thí nghiệm.}
{\True Trước khi chấp nhận kết luận, cần xem lại dữ liệu đo và các giả thiết, chẳng hạn việc công suất thất thoát có thực sự không đổi hay không.}
\loigiai{Đáp án: a) Đúng; b) Sai; c) Sai; d) Đúng. Xem lời giải và điều kiện vật lí tại hồ sơ PM01 trên Git.}
\end{ex}

